\documentclass[border=5pt]{standalone}

% --- Paquetes de Codificación y Lenguaje ---
\usepackage[utf8]{inputenc}
\usepackage[T1]{fontenc}
\usepackage[spanish, es-tabla]{babel}

% --- Paquetes Matemáticos y de Ingeniería ---
\usepackage{amsmath}
\usepackage{amsfonts}
\usepackage{amssymb}
\usepackage{siunitx} % Para magnitudes físicas (p.e. \SI{10}{\micro\farad})

% --- Configuración de Colores ---
\usepackage{xcolor}
\definecolor{black}{HTML}{000000} % Forzamos B/N estricto

% --- Núcleo de Diagramación: TikZ y CircuiTikZ ---
\usepackage{tikz}
\usepackage[american, cute inductors, american currents, american voltages]{circuitikz}

% --- Librerías de TikZ Especializadas ---
\usetikzlibrary{arrows.meta}    % Para puntas de flecha IEEE
\usetikzlibrary{positioning}   % Para posicionamiento relativo de nodos
\usetikzlibrary{calc}          % Para cálculos de coordenadas
\usetikzlibrary{shapes.geometric}

% --- Configuración de Estilo CircuiTikZ (Estándar IEEE) ---
\ctikzset{
	logic ports=ieee,                     % Compuertas según IEEE Std 91/91a
	resistors/scale=0.8,                  % Escalamiento profesional
	capacitors/scale=0.7,
	inductors/scale=0.7,
	diodes/scale=0.7,
	voltage/distance from node=0.4,       % Ajuste de etiquetas de voltaje
	bipoles/thickness=1                   % Grosor de línea para componentes
}

% --- Definición de Estilos Personalizados ---
\tikzset{
	label style/.style={font=\small\sffamily},
	connection/.style={circle, fill=black, inner sep=1.5pt} % Para nodos de conexión (junctions)
}

\begin{document}
	
\begin{circuitikz}[line width=0.8pt, american]
	% --- 1. ETAPA DE TRANSFORMACIÓN ---
	\node[transformer core] (T) at (0,0) {};
	\draw (T.A1) to[short, -o] ++(-1,0) node[left] {$V_{ac,+}$};
	\draw (T.A2) to[short, -o] ++(-1,0) node[left] {$V_{ac,-}$};
	
	% --- 2. ETAPA DE RECTIFICACIÓN (Puente Discreto) ---
	\coordinate (Bleft) at (2.5, 0);
	\coordinate (Bright) at (4.5, 0);
	\coordinate (Btop) at (3.5, 1.5);
	\coordinate (Bbot) at (3.5, -1.5);
	
	% Trazado de diodos (Ánodo en la primera coordenada, Cátodo en la segunda)
	\draw (Bleft) to[D] (Btop);
	\draw (Bright) to[D] (Btop);
	\draw (Bbot) to[D] (Bleft);
	\draw (Bbot) to[D] (Bright);
	
	% Nodos de unión físicos del puente rectificador
	\node[circle, fill, inner sep=1.2pt] at (Bleft) {};
	\node[circle, fill, inner sep=1.2pt] at (Bright) {};
	\node[circle, fill, inner sep=1.2pt] at (Btop) {};
	\node[circle, fill, inner sep=1.2pt] at (Bbot) {};
	
	% Enrutamiento ortogonal del transformador al puente
	\draw (T.B1) -| (Bleft);
	% T.B2 transita por y=-2.5 para evitar cruces ambiguos con el riel de tierra
	\draw (T.B2) -- ++(0.5,0) |- (4.5, -2.5) -- (Bright);
	
	% --- 3. RIELES DE ALIMENTACIÓN PRINCIPALES ---
	% Riel positivo de entrada al regulador (Sube a y=2.5)
	\draw (Btop) -- (3.5, 2.5) -- (8, 2.5);
	% Riel negativo de referencia (Baja a y=-3.5)
	\draw (Bbot) -- (3.5, -3.5) -- (14, -3.5);
	\draw (13.5, -3.5) node[ground] {};
	
	% --- 4. ETAPA DE FILTRADO PREVIO ---
	% Capacitor Electrolítico (Bulk) con polaridad '+' forzada en la parte superior
	\draw (5.5, 2.5) to[eC, l=$C_{bulk}$, *-*] (5.5, -3.5);
	% Capacitor cerámico para filtrado de alta frecuencia
	\draw (6.8, 2.5) to[C, l=$C_{in}$, *-*] (6.8, -3.5);
	
	% --- 5. ETAPA DE REGULACIÓN ---
	% Bloque del regulador lineal general
	\node[draw, minimum width=2cm, minimum height=1.2cm, fill=white, anchor=west] (REG) at (8, 2.5) {\small Regulador};
	
	% Riel positivo de salida del regulador
	\draw (REG.east) -- (14, 2.5);
	
	% Conexión del pin de referencia a tierra (Cálculo de intersección |-)
	\draw (REG.south) -- (REG.south |- 0, -3.5);
	\node[circle, fill, inner sep=1.2pt] at (REG.south |- 0, -3.5) {};
	
	% Diodo de protección (Flyback) OUT -> IN
	% Coordenadas relativas extraídas para mantener paralelismo estricto
	\draw (REG.east) ++(0.4, 0) coordinate (P_OUT);
	\draw (REG.west) ++(-0.4, 0) coordinate (P_IN);
	\draw (P_OUT) to[short, *-] ++(0, 1.5) to[D, l=$D_{prot}$] ++(-2.8, 0) to[short, -*] (P_IN);
	
	% --- 6. ETAPA DE FILTRADO POSTERIOR Y CARGA ---
	% Desacoplo de salida
	\draw (11.2, 2.5) to[C, l=$C_{out}$, *-*] (11.2, -3.5);
	% Estabilización de carga
	\draw (12.2, 2.5) to[eC, l=$C_{load}$, *-*] (12.2, -3.5);
	
	% Resistencia de purga / simulación de carga
	\draw (13.2, 2.5) to[R, l=$R_{load}$, *-*] (13.2, -3.5);
	
	% --- 7. PUERTOS DE SALIDA FINALES ---
	\draw (14, 2.5) to[short, -o] (14.5, 2.5) node[right] {$V_{out,+}$};
	\draw (14, -3.5) to[short, -o] (14.5, -3.5) node[right] {$V_{out,-}$};
	
\end{circuitikz}

\begin{circuitikz}[line width=0.8pt, american]
	% --- LIMITES DEL CHIP (SILICIO) ---
	% Se delimita el encapsulado del circuito integrado
	\draw[dashed, teal, line width=1.2pt, rounded corners=3mm] (0.8, -0.5) rectangle (11.5, 7.8);
	\node[teal, font=\bfseries\sffamily, right] at (1, 7.4) {Topología Interna del Regulador (Nivel de Silicio)};
	
	% --- PUERTOS EXTERNOS (3 Terminales) ---
	\draw (0, 7) node[left, font=\bfseries] {IN} to[short, o-*] (2, 7);
	\draw (12.5, 3.5) node[right, font=\bfseries] {OUT} to[short, o-*] (11, 3.5);
	\draw (6.5, -1.5) node[below, font=\bfseries] {GND / ADJ} to[short, o-*] (6.5, 0);
	
	% --- RIEL DE ENTRADA ---
	\draw (2, 7) -- (8, 7);
	
	% --- REFERENCIA DE VOLTAJE INTERNA ---
	% Polarización y diodo Zener para V_ref
	\draw (2, 7) to[R, l=$R_{bias}$] (2, 3.5) to[zzD, l=$V_{ref}$, -*] (2, 0);
	
	% --- AMPLIFICADOR DE ERROR ---
	% Se utiliza 'noinv input up' para mantener el flujo geométrico limpio hacia V_ref
	\node[op amp, noinv input up] (OA) at (5, 3) {};
	\node[above, yshift=-2pt] at (OA.center) {\tiny \sffamily Error};
	
	% Alimentación asimétrica del Op-Amp desde los rieles internos
	\draw (OA.up) -- (OA.up |- 0, 7) node[circle, fill, inner sep=1.2pt]{};
	\draw (OA.down) -- (OA.down |- 0, 0) node[circle, fill, inner sep=1.2pt]{};
	
	% Conexión del voltaje de referencia al pin No Inversor (+)
	\draw (2, 3.5) to[short, *-] (OA.+);
	
	% --- TRANSISTOR DE PASO (ELEMENTO EN SERIE) ---
	\node[npn] (Q1) at (8, 6) {};
	\node[above left=1pt] at (Q1.B) {$Q_{pass}$};
	
	% Toma de corriente del colector al riel principal de entrada
	\draw (Q1.C) -- (Q1.C |- 0, 7) node[circle, fill, inner sep=1.2pt]{};
	
	% Enrutamiento de la salida del Amp Error hacia la Base de Q_pass
	% Se define el nodo (NB) explícitamente para la etapa de limitación
	\draw (OA.out) -- (6, 3) |- (7, 6) node[circle, fill, inner sep=1.2pt] (NB) {} -- (Q1.B);
	
	% --- ETAPA DE PROTECCIÓN CONTRA SOBRECORRIENTE ---
	\node[npn] (Q2) at (10, 4.4) {};
	\node[right=2pt] at (Q2.B) {$Q_{limit}$};
	
	% Resistencia de sensado (R_sc) en serie con el emisor de potencia
	\draw (Q1.E) node[circle, fill, inner sep=1.2pt] (NE) {} to[R, l=$R_{sc}$, -*] (8, 3.5);
	
	% Sensado diferencial: La base lee aguas arriba de R_sc, el emisor aguas abajo
	\draw (NE) -- (9, 5.24) |- (Q2.B);
	\draw (Q2.E) -- (Q2.E |- 0, 3.5) node[circle, fill, inner sep=1.2pt]{};
	
	% El colector roba corriente de derivación (base drive) de Q_pass
	% El enrutamiento evita cruzar físicamente sobre el símbolo de Q_pass
	\draw (Q2.C) |- (10.5, 6.5) -| (NB);
	
	% --- RED DE RETROALIMENTACIÓN (DIVISOR INTERNO) ---
	\draw (8, 3.5) -- (11, 3.5);
	\draw (10.5, 3.5) to[R, l=$R_1$, *-*] (10.5, 1.75) to[R, l=$R_2$, -*] (10.5, 0);
	
	% Enrutamiento del nodo de muestra hacia el pin Inversor (-) del Amp Error
	\draw (10.5, 1.75) -- (3.2, 1.75) |- (OA.-);
	
	% --- RIEL COMÚN DE REFERENCIA ---
	\draw (2, 0) -- (10.5, 0);
	
\end{circuitikz}

\begin{circuitikz}[line width=0.8pt, american]
	
	% --- RIELES DE ALIMENTACIÓN ---
	\draw (1, 10) node[left, font=\bfseries] {$V_{CC}$} to[short, o-] (13, 10);
	\draw (1, 0) node[left, font=\bfseries] {$V_{EE}$} to[short, o-] (13, 0);
	
	% ==========================================
	% ETAPA 1: ENTRADA DIFERENCIAL Y ESPEJO
	% ==========================================
	% Par Diferencial NPN
	\node[npn] (Q1) at (4, 5.5) {}; \node[below right] at (Q1.B) {$Q_1$};
	\node[npn, xscale=-1] (Q2) at (7, 5.5) {}; \node[below left] at (Q2.B) {$Q_2$};
	
	% Entradas
	\draw (Q1.B) to[short, -o] ++(-1.5,0) node[left] {$V_{in-}$};
	\draw (Q2.B) to[short, -o] ++(1.5,0) node[right] {$V_{in+}$};
	
	% Espejo de Corriente de Carga Activa (PNP)
	\node[pnp] (Q3) at (4, 8) {}; \node[above right] at (Q3.B) {$Q_3$};
	\node[pnp, xscale=-1] (Q4) at (7, 8) {}; \node[above left] at (Q4.B) {$Q_4$};
	
	% Conexiones de Emisores a Vcc
	\draw (Q3.E) -- (4, 10) node[circle, fill, inner sep=1.2pt]{};
	\draw (Q4.E) -- (7, 10) node[circle, fill, inner sep=1.2pt]{};
	
	% Conexión de los Colectores (Q1-Q3 y Q2-Q4)
	\draw (Q1.C) -- (Q3.C);
	\draw (Q2.C) -- (Q4.C) node[midway, circle, fill, inner sep=1.2pt] (V_diff) {};
	
	% Enrutamiento de las bases del espejo (Evitando cruzar los símbolos)
	\draw (Q3.B) -- (3.2, 8) -- (3.2, 9) -- (7.8, 9) -- (7.8, 8) -- (Q4.B);
	
	% Conexión tipo diodo para Q3
	\draw (3.2, 8) |- (4, 6.8) node[circle, fill, inner sep=1.2pt] {};
	
	% Fuente de corriente de cola
	\draw (Q1.E) |- (5.5, 4.3) node[circle, fill, inner sep=1.2pt] {} -| (Q2.E);
	\draw (5.5, 4.3) to[I, l=$I_{tail}$] (5.5, 0) node[circle, fill, inner sep=1.2pt]{};
	
	% ==========================================
	% ETAPA 2: AMPLIFICACIÓN DE VOLTAJE (VAS)
	% ==========================================
	\node[npn] (Q5) at (9.5, 1.5) {}; \node[below right] at (Q5.B) {$Q_5$};
	\draw (Q5.E) -- (9.5, 0) node[circle, fill, inner sep=1.2pt]{};
	
	% Enrutamiento desde la salida diferencial hacia la base de Q5
	\draw (V_diff) -- (8.5, 6.75) |- (Q5.B);
	
	% ==========================================
	% RED DE POLARIZACIÓN Y ETAPA DE SALIDA
	% ==========================================
	% Fuente de corriente de polarización desde Vcc
	\draw (9.5, 10) node[circle, fill, inner sep=1.2pt]{} to[I, l=$I_{bias}$] (9.5, 8) node[circle, fill, inner sep=1.2pt] (NodeA) {};
	
	% Diodos de polarización para la clase AB
	\draw (NodeA) to[D, l=$D_1$] (9.5, 5.5) to[D, l=$D_2$, *-*] (9.5, 3) node[circle, fill, inner sep=1.2pt] (NodeB) {};
	
	% Conexión al colector del VAS
	\draw (NodeB) -- (Q5.C);
	
	% Capacitor de Compensación Miller (Cc)
	\draw (8.5, 1.5) node[circle, fill, inner sep=1.2pt]{} -- (8.5, 4) to[C, l=$C_c$, -*] (9.5, 4);
	
	% Transistores de Salida Push-Pull (Clase AB)
	\node[npn] (Q6) at (12, 8) {}; \node[above left] at (Q6.B) {$Q_6$};
	\node[pnp] (Q7) at (12, 3) {}; \node[below left] at (Q7.B) {$Q_7$};
	
	% Alimentación de los colectores de salida
	\draw (Q6.C) -- (12, 10) node[circle, fill, inner sep=1.2pt]{};
	\draw (Q7.C) -- (12, 0) node[circle, fill, inner sep=1.2pt]{};
	
	% Conexión de bases
	\draw (NodeA) -- (Q6.B);
	\draw (NodeB) -- (Q7.B);
	
	% Resistencias de degeneración de emisor y salida
	\draw (Q6.E) to[R, l=$R_{e1}$] (12, 5.5) node[circle, fill, inner sep=1.2pt] (OUT) {};
	\draw (Q7.E) to[R, l=$R_{e2}$] (OUT);
	
	% Puerto de Salida
	\draw (OUT) to[short, -o] (13.5, 5.5) node[right, font=\bfseries] {$V_{out}$};
	
	% ==========================================
	% DELIMITACIONES ARQUITECTÓNICAS (Anotaciones)
	% ==========================================
	\draw[dashed, gray] (2.5, 3.5) rectangle (8, 9.5);
	\node[gray, font=\sffamily\scriptsize, align=center] at (5.25, 9.2) {Etapa de\\Entrada Diferencial};
	
	\draw[dashed, gray] (8.2, 0.5) rectangle (10.5, 9.5);
	\node[gray, font=\sffamily\scriptsize, align=center, rotate=90] at (10.2, 5) {Etapa VAS y Polarización};
	
	\draw[dashed, gray] (10.8, 1.5) rectangle (12.8, 9.5);
	\node[gray, font=\sffamily\scriptsize, align=center, rotate=90] at (12.5, 5.5) {Etapa de Salida Clase AB};
	
\end{circuitikz}

\begin{circuitikz}[line width=0.8pt, american]
	
	% =======================================================
	% MACRO CORREGIDA: OSCILOSCOPIO MATEMÁTICO VIRTUAL
	% =======================================================
	\def\oscope#1#2#3#4#5{
		% #1: Posición X, #2: Posición Y, #3: Etiqueta, #4: Nodo de sondeo, #5: Ecuación/Path
		\begin{scope}[shift={(#1,#2)}]
			% Chasis del osciloscopio
			\draw[fill=white, thick, rounded corners=2pt] (-0.2, -1.2) rectangle (3.2, 1.4);
			% Rejilla de la pantalla (Graticule)
			\draw[step=0.25, gray!30, thin] (0, -1) grid (3, 1);
			% Ejes X e Y
			\draw[gray!80, thin] (0, 0) -- (3, 0); 
			\draw[gray!80, thin] (0, -1) -- (0, 1);
			% Etiqueta del instrumento
			\node[above, font=\sffamily\tiny\bfseries] at (1.5, 1.05) {#3};
			
			% Trazado de la onda
			\begin{scope}[shift={(0,0)}]
				\draw[thick, black] #5;
			\end{scope}
			
			% Punto de anclaje de la sonda en el instrumento
			\coordinate (osc_bot) at (1.5, -1.2);
		\end{scope}
		
		% CONEXIÓN INFALIBLE: Línea punteada simple + Nodo circular explícito
		\draw[densely dashed, gray, thick] (#4) -- (osc_bot);
		\fill[gray] (#4) circle (1.5pt);
	}
	
	% =======================================================
	% TOPOLOGÍA DEL CIRCUITO
	% =======================================================
	
	% --- 1. ETAPA DE TRANSFORMACIÓN ---
	\node[transformer core] (T) at (0,0) {};
	% Nodos de entrada
	\draw (T.A1) to[short, -o] ++(-1.5,0) coordinate (N1) node[left] {$V_{ac,+}$};
	\draw (T.A2) to[short, -o] ++(-1.5,0) node[left] {$V_{ac,-}$};
	
	% Nodo del secundario
	\coordinate (N2) at (T.B1);
	
	% --- 2. ETAPA DE RECTIFICACIÓN ---
	\coordinate (Bleft) at (2.5, 0);
	\coordinate (Bright) at (4.5, 0);
	\coordinate (Btop) at (3.5, 1.5);
	\coordinate (Bbot) at (3.5, -1.5);
	
	\draw (Bleft) to[D] (Btop);
	\draw (Bright) to[D] (Btop);
	\draw (Bbot) to[D] (Bleft);
	\draw (Bbot) to[D] (Bright);
	
	\node[circle, fill, inner sep=1.2pt] at (Bleft) {};
	\node[circle, fill, inner sep=1.2pt] at (Bright) {};
	\node[circle, fill, inner sep=1.2pt] (N3) at (Btop) {}; % Nodo Rectificado
	\node[circle, fill, inner sep=1.2pt] at (Bbot) {};
	
	\draw (T.B1) -| (Bleft);
	\draw (T.B2) -- ++(0.5,0) |- (4.5, -2.5) -- (Bright);
	
	% Rieles principales
	\draw (Btop) -- (3.5, 2.5) -- (8, 2.5);
	\draw (Bbot) -- (3.5, -3.5) -- (14, -3.5);
	\draw (13.5, -3.5) node[ground] {};
	
	% --- 3. ETAPA DE FILTRADO ---
	\draw (5.5, 2.5) node[circle, fill, inner sep=1.2pt] (N4) {} to[eC, l=$C_{bulk}$, *-*] (5.5, -3.5);
	
	% --- 4. ETAPA DE REGULACIÓN ---
	\node[draw, minimum width=2cm, minimum height=1.2cm, fill=white, anchor=west] (REG) at (8, 2.5) {\small Regulador};
	\draw (REG.south) -- (REG.south |- 0, -3.5);
	\node[circle, fill, inner sep=1.2pt] at (REG.south |- 0, -3.5) {};
	
	% --- 5. ETAPA DE SALIDA ---
	\draw (REG.east) -- (14, 2.5) coordinate (N5);
	\node[circle, fill, inner sep=1.2pt] at (11.2, 2.5) {};
	\draw (11.2, 2.5) to[C, l=$C_{out}$, *-*] (11.2, -3.5);
	\draw (N5) to[short, -o] (14.5, 2.5) node[right] {$V_{out,+}$};
	\draw (14, -3.5) to[short, -o] (14.5, -3.5) node[right] {$V_{out,-}$};
	
	% =======================================================
	% INYECCIÓN DE LOS OSCILOSCOPIOS (ANÁLISIS DE SEÑAL)
	% =======================================================
	
	\oscope{-2}{6}{1. Entrada Red AC}{N1}{ plot[domain=0:2.8, samples=100] (\x, {0.9*sin(2.5*\x r)}) }
	
	\oscope{1.5}{6}{2. Secundario}{N2}{ plot[domain=0:2.8, samples=100] (\x, {0.45*sin(2.5*\x r)}) }
	
	\oscope{5}{6}{3. Rectificado}{N3}{ plot[domain=0:2.8, samples=150] (\x, {0.45*abs(sin(2.5*\x r))}) }
	
	\oscope{8.5}{6}{4. Offset y Rizado}{N4}{ (0, 0.25) -- (0.63, 0.45) -- (1.7, 0.25) -- (1.88, 0.45) -- (2.8, 0.25) }
	
	\oscope{12}{6}{5. Regulador DC}{N5}{ (0, 0.15) -- (2.8, 0.15) }
	
\end{circuitikz}
\end{document}